\subsection{Remaining-Sequence Admission}
\label{sec:admission}

A capture's timeout window continues to shrink while it waits for the draft worker and executes earlier stages.
\emph{Remaining-Sequence Admission} therefore reads the live budget immediately before each optional stage.
The decision accounts for the entire remaining sequence, including mandatory encoding and saving, so it considers whether the draft image can still be produced after admitting the stage.

\smallskip\noindent\textbf{Live-budget admission.}
At admission point \(j\), let \(T_{i,j}\) be the time remaining until capture \(i\)'s Capture Timeout deadline.
Let \(U(\mathcal K_{i,j})\) be an empirical high-side estimate of its remaining processing time, obtained by correcting \(\hat P\) with observed prediction errors as described below.
An optional stage in an enabled group is admitted when
\begin{equation}
  U(\mathcal{K}_{i,j})\leq T_{i,j}
  \quad\text{or}\quad
  \hat P(\mathcal{K}_{i,j})=0.
  \label{eq:admission}
\end{equation}
A zero point estimate permits execution to obtain the initial duration observations needed by the model.
This cold-start decision has no learned timing evidence; any Capture Timeout observed during its validation remains a release blocker.

Once an \(M\) or \(S\) stage is rejected, the controller skips the optional stages in that group and removes them from subsequent remaining sequences until the queue drains.
The two groups are tracked independently, so rejecting one does not disable the other.
Keeping a group disabled until the drain prevents optional effects from alternating across consecutive captures.

\smallskip\noindent\textbf{Residual-calibrated upper estimate.}
The point estimate \(\hat P\) can understate processing time when runtime conditions or image content cause a stage to run longer than predicted.
To account for these overruns, the controller learns a multiplicative correction from completed sequences.
For each admission decision, the controller records the per-key predictions \(\hat p(k)\) for the corresponding remaining sequence \(\mathcal K\).
When the draft sequence completes, the controller compares those predictions with the realized durations \(p(k)\).
A recorded \(\mathcal K\) contributes a residual factor only if \(\hat P(\mathcal K)>0\) and \(p(k)>0\) for all \(k\in\mathcal K\):
\begin{equation}
  \phi(\mathcal{K})
  =\frac{\sum_{k\in\mathcal{K}}\max\{\hat p(k),p(k)\}}
         {\sum_{k\in\mathcal{K}}\hat p(k)}.
  \label{eq:residual-factor}
\end{equation}
A skipped stage therefore excludes the recorded sequences containing it, while fully measured remaining sequences can still contribute.
The per-key maximum preserves each stage's overrun even when another stage finishes faster than predicted, ensuring \(\phi(\mathcal{K})\geq1\).
Normalizing by the predicted duration makes the factor comparable across sequences.

When a completion yields factors, the controller multiplies the existing weights in all sequence histories and the cross-sequence pool by \(0.90\).
It then inserts each new factor with unit weight into its sequence history and the pool.
To calibrate a remaining sequence, it selects the weighted \(\tau\)-quantile from that sequence's history.
If the history is empty, it uses the pool; if both are empty, it sets \(\phi_\tau(\mathcal K)=1\).
It determines \(\tau\) from the Kish effective sample size \(n_{\mathrm{eff}}\) of the current weights \(w_l\) in the selected history~\cite{kish-survey-sampling}:
\[
  n_{\mathrm{eff}}
  =\frac{(\sum_l w_l)^2}{\sum_l w_l^2},
  \qquad
  \tau=\frac{n_{\mathrm{eff}}}{n_{\mathrm{eff}}+1}.
\]
The selector is motivated by the expected CDF position \(n/(n+1)\) of the maximum of \(n\) i.i.d. continuous samples~\cite{david-order-statistics}, with \(n_{\mathrm{eff}}\) substituting for the sample count.
As the effective sample size grows, the selected quantile moves further into the observed upper tail.
Recency weighting lets old observations lose influence without requiring a separate target percentile for each device or stage.
Applying the selected factor gives
\begin{equation}
  U(\mathcal{K})
  =\hat P(\mathcal{K})\,\phi_{\tau}(\mathcal{K}).
  \label{eq:upper-estimate}
\end{equation}
The resulting \(U\) is empirical; it provides neither a worst-case time bound nor a statistical coverage guarantee.

\smallskip\noindent\textbf{Watchdog fallback.}
Even an admitted stage can slow down after the decision.
The watchdog therefore limits how long draft completion waits for that stage before starting fallback.
It sets the waiting interval to
\begin{equation}
  W_{i,j}=\max\{0,T_{i,j}-U(\mathcal{E}_i)\},
  \label{eq:watchdog}
\end{equation}
where \(\mathcal{E}_i\) contains only mandatory processing from encoding through draft image saving.
Subtracting \(U(\mathcal E_i)\) reserves its estimated duration within the remaining budget.
If the watchdog expires, the framework proceeds to encoding and saving from a preserved copy of the original input.
Already queued draft sequences follow the same fallback until the queue drains.
Expiration does not stop native processing already in progress; Section~\ref{sec:implementation} explains how fallback proceeds while that processing finishes.
